\subsection[gretl.Matrix class]{\classpymatrix}\label{sec:classpymatrix}
A \gretlcmd{matrix} in \gretl behaves much like a matrix in other
mathematical software (e.g. MATLAB or Octave). Matrices have no
ID numbers and are referenced by name. They can also be used
independently of any dataset. Unlike the native \gretl clients,
which use 1-based indexing, \gretlmatrix indexing is 0-based:
the top-left element of matrix \pythcode{A} is \pythcode{A[0,0]}.

The \classpymatrix has one important advantage over
\pythcode{numpy.ndarray}: it may have row and/or column names
attached, which are displayed when the matrix is printed.


\subsubsection{\classpymatrixb constructors}
The \classpymatrix provides the following constructors:

\gretlfpycmd{Matrix()}
\estdesc{Default constructor.}
\begin{lstlisting}[style=pythonstylefoot]
~gretl.Matrix -> gretl.Matrix
\end{lstlisting}
\maptof{\{\}}{matrix}
Example:\newline
\pythcodef{m = ~gretl.Matrix()}
\generaldesc{Constructs an empty \pythcode{gretl.Matrix} object.}

\gretlfpycmd{Matrix()}
\estdesc{Creates a $rows\times cols$ real or complex
	\pythcode{gretl.Matrix}, optionally filled with \textit{fill}.}
\begin{lstlisting}[style=pythonstylefoot]
~gretl.Matrix(rows: int, cols: int[, fill: float = 0.0, type: str = 'real']) -> gretl.Matrix
\end{lstlisting}
\maptof{ones(.rows, .cols) .* fill}{matrix}
Example:\newline
\pythcodef{m = ~gretl.Matrix(3, 2, 3.4)}\\
\pythcodef{m = ~gretl.Matrix(3, 2, -1.4, 'complex')}
\generaldesc{Constructs a $rows\times cols$ \pythcode{gretl.Matrix},
	optionally filled with \textit{fill}. For
	\pythcode{type = 'complex'}, the user gets a complex matrix
	with real parts set to \textit{fill}.}

\gretlfpycmd{Matrix()}
\estdesc{Creates a $rows\times cols$ random \pythcode{gretl.Matrix}.}
\begin{lstlisting}[style=pythonstylefoot]
~gretl.Matrix(rows: int, cols: int[, dist: str = 'normal']) -> gretl.Matrix
\end{lstlisting}
\maptof{mnormal(.rows, .cols)}{matrix}
\maptof{muniform(.rows, .cols)}{matrix}
Example:\newline
\pythcodef{m = ~gretl.Matrix(3, 2, 'normal')}\\
\pythcodef{m = ~gretl.Matrix(3, 4, 'U')}
\generaldesc{Constructs a $rows\times cols$ random \pythcode{gretl.Matrix}.
	The pseudo-random values are drawn from one of the following distributions:
	\begin{itemize}
		\item Standard Normal distribution:
		\pythcode{dist = 'normal' | 'n' | 'N' | 'z' | 'Z'},
		\item Uniform distribution:
		\pythcode{dist = 'uniform' | 'u' | 'U'}.
	\end{itemize}}

\gretlfpycmd{Matrix()}
\estdesc{Creates a row vector containing a sequence of values
	from \textit{a} to \textit{b}.}
\begin{lstlisting}[style=pythonstylefoot]
~gretl.Matrix(a: float, b: float[, k: float = 1.0]) -> gretl.Matrix
\end{lstlisting}
\maptof{seq(a, b, k)}{matrix}
Example:\newline
\pythcodef{m = ~gretl.Matrix(1.0, 5.0)}\\
\pythcodef{m = ~gretl.Matrix(10.0, 4.0, 2.0)}
\generaldesc{Given two arguments, returns a row vector
	containing values from \textit{a} to \textit{b} with
	a step size of~1, or~$-1$ if \textit{a} is greater
	than \textit{b}. If the third argument is given,
	returns a row vector containing a sequence of values
	starting with \textit{a} and incremented (or
	decremented, if \textit{a} is greater than
	\textit{b}) by \textit{k} at each step. The final
	value is the largest member of the sequence that is
	less than or equal to \textit{b} (or, mutatis
	mutandis, greater than or equal to \textit{b} if
	\textit{a} is greater than \textit{b}). The argument
	\textit{k} must be positive.}

\gretlfpycmd{Matrix()}
\estdesc{Creates an identity matrix.}
\begin{lstlisting}[style=pythonstylefoot]
~gretl.Matrix(n: int) -> gretl.Matrix
\end{lstlisting}
\maptof{I(n)}{matrix}
Example:\newline
\pythcodef{m = ~gretl.Matrix(3)}
\generaldesc{Constructs an $n\times n$ \pythcode{gretl.Matrix}
	with ones on the main diagonal and zeros elsewhere.}

\gretlfpycmd{Matrix()}
\estdesc{Creates a \pythcode{gretl.Matrix} from a \pythcode{numpy.ndarray}.}
\begin{lstlisting}[style=pythonstylefoot]
~gretl.Matrix(array: numpy.ndarray) -> gretl.Matrix
\end{lstlisting}
Example:\newline
\pythcodef{v = ~gretl.Matrix(np.array([1, 2, 3, 4]))}\\
\pythcodef{m = ~gretl.Matrix(np.array([[1, 2], [3, 4]]))}
\generaldesc{Constructs a \pythcode{gretl.Matrix} as a deep
	copy of an existing \pythcode{numpy.ndarray}. If the input is
	a one-dimensional NumPy array, it is interpreted as a column
	vector. Thus, in the above example, \textit{v} is a column vector.}

\gretlfpycmd{Matrix()}
\estdesc{Creates a \pythcode{gretl.Matrix} from a Python \pythcode{Sequence}.}
\begin{lstlisting}[style=pythonstylefoot]
~gretl.Matrix(seq: Sequence) -> gretl.Matrix
\end{lstlisting}
Example:\newline
\pythcodef{v = ~gretl.Matrix([1, 2, 3, 4])}\\
\pythcodef{m = ~gretl.Matrix([[1, 2], [3, 4]])}
\generaldesc{Constructs a \pythcode{gretl.Matrix} from a
	Python \pythcode{Sequence}. The input \pythcode{seq} must be a
	Python sequence (e.g., a \pythcode{list} or \pythcode{tuple}).
	A one-dimensional sequence is interpreted as a row vector,
	while a nested sequence is interpreted as a matrix whose rows
	correspond to the inner sequences. All rows of a nested sequence
	must have the same length, and all elements must be numeric.
	Thus, in the above example, \textit{v} is a row vector,
	while \textit{m} is a matrix whose rows correspond to the
	inner lists.}


\subsubsection{\classpymatrixb operators}
The \classpymatrix class provides the following operators:

\gretlfpycmd{[]}
\estdesc{Operator for accessing or assigning \pythcode{gretl.Matrix} elements.}
\begin{lstlisting}[style=pythonstylefoot]
v[i] -> float | complex
m[i,j] -> float | complex
m[i:j] -> gretl.Matrix
m[i:j,k:l] -> gretl.Matrix
v[i] = float | complex
m[i,j] = float | complex
\end{lstlisting}
Example:\\
\pythcodef{v = m[0]}\\
\pythcodef{f = m[0,0]}\\
\pythcodef{s = m[:,-1]}\\
\pythcodef{v[0] = 2 - .5j}\\
\pythcodef{m[0,0] = 10 + 1j}
\generaldesc{The operator retrieves or assigns values of an
	existing \pythcode{gretl.Matrix} using Python indexing syntax.
	Scalar indexing returns a numeric value, while slicing returns a
	new \pythcode{gretl.Matrix}. Both vector and matrix scalar elements
	can be assigned. Slice assignment is not supported; slices are
	available for retrieving data only.}


\subsubsection{\classpymatrixb attributes}
The \classpymatrix provides the following public attributes:

\gretlfpycmd{cols}
\estdesc{Returns the number of columns of the matrix.}
\begin{lstlisting}[style=pythonstylefoot]
.cols -> int
\end{lstlisting}

\gretlfpycmd{is\_complex}
\estdesc{Returns \pythcode{True} if the matrix is complex-valued.}
\begin{lstlisting}[style=pythonstylefoot]
.is_complex -> bool
\end{lstlisting}

\gretlfpycmd{is\_vector}
\estdesc{Returns \pythcode{True} if the matrix is a vector
	(that is, a one-dimensional matrix).}
\begin{lstlisting}[style=pythonstylefoot]
.is_vector -> bool
\end{lstlisting}

\gretlfpycmd{numpy\_view}
\estdesc{Returns a \pythcode{numpy.ndarray} view of the \pythcode{gretl.Matrix}.}
\begin{lstlisting}[style=pythonstylefoot]
.numpy_view -> numpy.ndarray
\end{lstlisting}
\generaldesc{This property returns a \pythcode{numpy.ndarray}
	view sharing the underlying storage of the
	\pythcode{gretl.Matrix}. No data copying is performed,
	so the operation has negligible overhead. The returned
	view is writable, but its shape cannot be changed.}

\gretlfpycmd{pandas\_view}
\estdesc{Returns a \pythcode{pandas.DataFrame} constructed from
	the \pythcode{gretl.Matrix} using a NumPy view.}
\begin{lstlisting}[style=pythonstylefoot]
.pandas_view -> pandas.core.frame.DataFrame
\end{lstlisting}
\generaldesc{This property returns a \pythcode{pandas.DataFrame}
	constructed from a zero-copy NumPy view of the underlying
	\pythcode{gretl.Matrix} data. No explicit data copying is
	performed. However, since \pythcode{pandas.DataFrame} does
	not guarantee full view semantics, modifications to the
	returned object may not always propagate back to the
	original matrix.}

\gretlfpycmd{rows}
\estdesc{Returns the number of rows of the matrix.}
\begin{lstlisting}[style=pythonstylefoot]
.rows -> int
\end{lstlisting}


\subsubsection{\classpymatrixb static methods}
The \classpymatrix class provides the following static methods:

\gretlfpycmd{Matrix.diag()}
\estdesc{Creates a square \pythcode{gretl.Matrix} with
	the main diagonal set to \textit{seq}.}
\begin{lstlisting}[style=pythonstylefoot]
~gretl.Matrix.diag(seq: Sequence) -> gretl.Matrix
\end{lstlisting}
Example:\newline
\pythcodef{m = ~gretl.Matrix.diag([1, 2, 3])}
\generaldesc{Constructs an $n\times n$ \pythcode{gretl.Matrix},
	where \textit{n} is the length of \textit{seq}, with
	the main diagonal set to the elements of \textit{seq}.
	All off-diagonal elements are set to zero.}

\gretlfpycmd{Matrix.eye()}
\estdesc{Creates an identity matrix.}
\begin{lstlisting}[style=pythonstylefoot]
~gretl.Matrix.eye(n: int[, m: int]) -> gretl.Matrix
\end{lstlisting}
\maptof{I(n, m)}{matrix}
Example:\newline
\pythcodef{m = ~gretl.Matrix.I(3)}\\
\pythcodef{m = ~gretl.Matrix.I(3, 4)}
\generaldesc{Constructs an $n\times n$ or $n\times m$
	\pythcode{gretl.Matrix} with ones on the main diagonal
	and zeros elsewhere.}

\gretlfpycmd{Matrix.ones()}
\estdesc{Creates a matrix filled with~1.}
\begin{lstlisting}[style=pythonstylefoot]
~gretl.Matrix.ones(rows: int, cols: int) -> gretl.Matrix
\end{lstlisting}
\maptof{ones(.rows, .cols)}{matrix}
Example:\newline
\pythcodef{m = ~gretl.Matrix.ones(3, 4)}
\generaldesc{Constructs a \textit{rows}$\times$\textit{cols}
	\pythcode{gretl.Matrix} with all elements set to one.}

\gretlfpycmd{Matrix.zeros()}
\estdesc{Creates a matrix filled with~0.}
\begin{lstlisting}[style=pythonstylefoot]
~gretl.Matrix.zeros(rows: int, cols: int) -> gretl.Matrix
\end{lstlisting}
\maptof{zeros(.rows, .cols)}{matrix}
Example:\newline
\pythcodef{m = ~gretl.Matrix.zeros(3, 4)}
\generaldesc{Constructs a \textit{rows}$\times$\textit{cols}
	\pythcode{gretl.Matrix} with all elements set to zero.}


\subsubsection{\classpymatrixb public methods}
The \classpymatrix class provides the following public methods:

\gretlfpycmd{boxplot()}\label{met:gretlmatrix::boxplot}
\estdesc{Generates boxplots for matrix columns using \lgretl and Matplotlib.}
\begin{lstlisting}[style=pythonstylefoot]
.boxplot(cols: str | list = ""[, OPTS]) -> matplotlib.figure.Figure .| None
\end{lstlisting}
\mapto[None]{boxplot --matrix=M}
\mapto{boxplot cols --matrix=M}
\begin{itemize}\setlength{\itemsep}{0.3em}\setlength{\parskip}{0pt}\setlength{\parsep}{0pt}
	\item[]\maptoopt{notches=True}{--notches}{show a 90 percent confidence interval for the median}
	\item[]\maptoopt{output="filename"}{--output=filename}{send output to the specified file}
\end{itemize}
Examples:\newline
\pythcodef{m = gretl.Matrix(100, 4, "normal")}\\
\pythcodef{fig = m.boxplot()}\\
\pythcodef{fig = m.boxplot("1 2 3")}\\
\pythcodef{fig = m.boxplot([1, 3])}\\
\pythcodef{fig = m.boxplot("1 2", notches=True)}\\
\pythcodef{m.boxplot(output="boxplot.png")}
\generaldesc{Generates boxplots representing the distribution of one
	or more matrix columns. The optional \texttt{cols} argument
	specifies the columns to be included, using 1-based numbering.
	It may be supplied as a space-separated string or as a Python
	list containing integers or their string representations.
	If omitted, all columns are included.
	
	Each box extends from the first quartile to the third quartile,
	enclosing the middle 50 percent of the observations. The
	horizontal line inside the box indicates the median, while
	a plus sign denotes the mean.
	
	The whiskers extend up to 1.5 times the interquartile range
	beyond the first and third quartiles. Observations outside
	this range are treated as outliers and displayed as individual
	points.
	
	Setting \textit{notches=True} adds a bootstrap-based 90 percent
	confidence interval for the median, represented by dashed
	horizontal lines.
	
	The \gretlcmd{--matrix} option is supplied automatically.
	The \textit{output} option controls graphical output: a
	Matplotlib figure is generated by default, while a filename
	directs the graph to the specified output file.
	
	Plot generation is performed by \lgretl, while the resulting
	\gnuplot script is translated into Python code by \gptranslate
	and rendered using Matplotlib. The method returns a Matplotlib
	Figure object, or \pythcode{None} when output is directed to
	a file.}

\gretlfpycmd{cnameget()}
\estdesc{Returns the column names of the \classpymatrix as a \pythcode{list}.}
\begin{lstlisting}[style=pythonstylefoot]
.cnameget(col: int = None) -> list | str
\end{lstlisting}
\generaldesc{This method returns the column names of the
	\pythcode{gretl.Matrix} object as a \pythcode{list}. If the optional
	\textit{col} argument is given, the name of the specified
	column is returned as a \pythcode{str}. If the matrix has no
	column names attached, an empty \pythcode{list} is returned.}

\gretlfpycmd{cnameset()}
\estdesc{Sets the column names of the \classpymatrix.}
\begin{lstlisting}[style=pythonstylefoot]
.cnameset(names: list)
\end{lstlisting}
\generaldesc{This method sets the column names of the
	\pythcode{gretl.Matrix} object using the provided
	\pythcode{list}. The number of elements in \textit{names}
	must equal the number of rows of the matrix.}

\gretlfpycmd{copy()}
\estdesc{Returns a deep copy of the \classpymatrix object.}
\begin{lstlisting}[style=pythonstylefoot]
.copy() -> gretl.Matrix
\end{lstlisting}
\mapto[none]{mat\_copy = mat}
\generaldesc{This method returns a new, fully independent
	\pythcode{gretl.Matrix} object containing the same
	elements as the original matrix. Any modifications made
	to the copy do not affect the original object.
	This method is equivalent to:
	\pythcodef{import copy; mat\_copy = copy.copy(mat)}.}

\gretlfpycmd{plot()}\label{met:gretlmatrix::plot}
\estdesc{Generates X--Y and time-series plots of matrix columns
	using \lgretl and Matplotlib.}
\begin{lstlisting}[style=pythonstylefoot]
.plot(cols: str | list = ""[, OPTS]) -> matplotlib.figure.Figure .| None
\end{lstlisting}
\mapto[None]{gnuplot --matrix=M}
\mapto[None]{gnuplot cols --matrix=M}
\mapto{gnuplot cols --matrix=M --time-series}
\begin{itemize}\setlength{\itemsep}{0.3em}\setlength{\parskip}{0pt}\setlength{\parsep}{0pt}
	\item[]\maptoopt{with\_lines=True}{--with-lines}{draw lines instead of points}
	\item[]\maptoopt{with\_lp=True}{--with-lp}{draw lines and points}
	\item[]\maptoopt{with\_impulses=True}{--with-impulses}{draw vertical impulses}
	\item[]\maptoopt{with\_boxes=True}{--with-boxes}{draw histogram-style boxes}
	\item[]\maptoopt{with\_steps=True}{--with-steps}{draw step graphs}
	\item[]\maptoopt{time\_series=True}{--time-series}{plot against time}
	\item[]\maptoopt{single\_yaxis=True}{--single-yaxis}{force a single Y axis}
	\item[]\maptoopt{y2axis=colspec}{--y2axis=colspec}{specify columns for the secondary Y axis}
	\item[]\maptoopt{ylogscale=base}{--ylogscale[=base]}{use a logarithmic Y axis}
	\item[]\maptoopt{fit=fitspec}{--fit=fitspec}{specify a fitted line}
	\item[]\maptoopt{font=fontspec}{--font=fontspec}{specify the plot font}
	\item[]\maptoopt{output="filename"}{--output=filename}{send output to the specified file}
\end{itemize}
Examples:\newline
\pythcodef{m = gretl.Matrix(100, 4, "normal")}\\
\pythcodef{fig = m.plot()}\\
\pythcodef{fig = m.plot("2 1")}\\
\pythcodef{fig = m.plot([2, 1], with\_lines=True)}\\
\pythcodef{fig = m.plot("1 2", time\_series=True, with\_lines=True)}\\
\pythcodef{fig = m.plot("2 1", fit="quadratic")}\\
\pythcodef{m.plot(output="matrix_plot.png")}
\generaldesc{Generates graphical representations of matrix columns
	using the native \lgretl plotting functionality. The optional
	\texttt{cols} argument specifies the columns to be plotted,
	using 1-based numbering. It may be supplied as a
	space-separated string or as a Python list containing
	integers or their string representations.
	
	For ordinary X--Y plots, the last selected column is
	interpreted as the X variable, while preceding columns
	are plotted on the Y axis. If \textit{cols} is omitted,
	the last column of the matrix is used as the X variable
	and the preceding columns are used as Y variables.
	When \textit{time\_series=True}, the selected columns
	are plotted against the observation index.
	
	By default, observations are represented by points. Various
	plotting styles are available, including lines, impulses,
	boxes, and steps. For multiple Y columns, a secondary axis
	may be selected automatically or explicitly using
	\textit{y2axis}.
	
	The \textit{fit} option allows fitted curves to be added
	to bivariate plots. As with the corresponding \lgretl
	command, plotting-style options may accept an optional
	column specification to restrict the style to selected
	series.
	
	The \gretlcmd{--matrix} option is supplied automatically,
	so the matrix does not need to be registered explicitly
	in the \lgretl namespace.
	
	The \textit{output} option controls graphical output.
	By default, a Matplotlib figure is generated. Setting
	\pythcode{output="display"} explicitly requests graphical
	output, while a filename directs the graph to the
	specified file.
	
	Plot generation is performed by \lgretl, while the resulting
	\gnuplot script is translated into Python code by \gptranslate
	and rendered using Matplotlib. The method returns a Matplotlib
	Figure object, or \pythcode{None} when output is directed
	to a file.}

\gretlfpycmd{rnameget()}
\estdesc{Returns the row names of the \classpymatrix as a \pythcode{list}.}
\begin{lstlisting}[style=pythonstylefoot]
.rnameget(row: int = None) -> list | str
\end{lstlisting}
\generaldesc{This method returns the row names of the
	\pythcode{gretl.Matrix} object as a \pythcode{list}. If the optional
	\textit{row} argument is given, the name of the specified
	row is returned as a \pythcode{str}. If the matrix has no
	row names attached, an empty \pythcode{list} is returned.}

\gretlfpycmd{rnameset()}
\estdesc{Sets the row names of the \classpymatrix.}
\begin{lstlisting}[style=pythonstylefoot]
.rnameset(names: list)
\end{lstlisting}
\generaldesc{This method sets the row names of the
	\pythcode{gretl.Matrix} object using the provided
	\pythcode{list}. The number of elements in \textit{names}
	must equal the number of rows of the matrix.}

\gretlfpycmd{scatters()}\label{met:gretlmatrix::scatters}
\estdesc{Generate multiple pairwise scatterplots of matrix columns
	on a single canvas.}
\begin{lstlisting}[style=pythonstylefoot]
.scatters(cols: str | list[, OPTS]) -> matplotlib.figure.Figure .| None
\end{lstlisting}
\mapto{scatters ycols ; xcols --matrix=M}
\begin{itemize}\setlength{\itemsep}{0.3em}\setlength{\parskip}{0pt}\setlength{\parsep}{0pt}
	\item[]\maptoopt{with_lines=True}{--with-lines}{draw line graphs instead of scatterplots}
	\item[]\maptoopt{output="filename"}{--output=filename}{send output to the specified file}
\end{itemize}
Examples:\newline
\pythcodef{m = gretl.Matrix(100, 4, "normal")}\\
\pythcodef{fig = m.scatters("1 ; 2 3 4")}\\
\pythcodef{fig = m.scatters("1 2 3 ; 4")}\\
\pythcodef{fig = m.scatters("1 2 3 ; 4", with\_lines=True)}\\
\pythcodef{m.scatters("1 ; 2 3 4", output="scatters.png")}
\generaldesc{Generates multiple pairwise graphs of one matrix
	column against several other columns. The \texttt{cols}
	argument contains two groups of column numbers separated
	by a semicolon. Column numbering is 1-based.
	
	Either the first group contains a single column to be
	plotted on the y-axis against all columns in the second
	group, or the second group contains a single column to
	be plotted on the x-axis against all columns in the first
	group.
	
	For example, \pythcode{m.scatters("1 ; 2 3 4")} generates
	three graphs, plotting the first column against columns
	2, 3, and 4, respectively. Conversely,
	\pythcode{m.scatters("1 2 3 ; 4")} plots the first three
	columns against the fourth column.
	
	The semicolon separates the two groups of columns.
	In the current implementation, a string argument should
	be used to specify this separator.
	
	By default, observations are represented as points. Setting
	\texttt{with\_lines=True} produces line graphs instead.
	
	A maximum of 16 plots can be generated on a single canvas;
	any additional columns are ignored. Such collections of
	scatterplots are useful for exploratory data analysis,
	particularly for examining relationships between matrix
	columns.
	
	The \gretlcmd{--matrix} option is supplied automatically.
	The \textit{output} option controls graphical output, which
	is generated by \lgretl, translated by \gptranslate, and
	rendered using Matplotlib.
	
	The method returns a Matplotlib Figure object, or
	\pythcode{None} when output is directed to a file.
	See also \pythcode{::plot()}
	(\ref{met:gretlmatrix::plot}) for individual X--Y
	and time-series plots.}

\gretlfpycmd{summary()}\label{met:gretlmatrix::summary}
\estdesc{Computes summary statistics for the columns of the matrix.}
\begin{lstlisting}[style=pythonstylefoot]
.summary([OPTS]) -> gretl.Matrix
\end{lstlisting}
\mapto{summary --matrix=matname}
\begin{itemize}\setlength{\itemsep}{0.3em}\setlength{\parskip}{0pt}\setlength{\parsep}{0pt}
	\item[]\maptoopt{simple=True}{--simple}{basic statistics only}
	\item[]\maptoopt{quiet=True}{--quiet}{suppress printed output}
\end{itemize}
\generaldesc{Computes descriptive statistics for the columns of the
	\classpymatrix and returns the resulting table as a \classpymatrix.}

\gretlfpycmd{to\_numpy()}
\estdesc{Returns a deep copy of the \classpymatrix as a NumPy array.}
\begin{lstlisting}[style=pythonstylefoot]
.to_numpy() -> numpy.ndarray
\end{lstlisting}
\generaldesc{This method returns a new, independent
	\pythcode{numpy.ndarray} containing the same data as the
	original \pythcode{gretl.Matrix}. Changes made to the
	returned array do not affect the original matrix.}

\gretlfpycmd{to\_pandas()}
\estdesc{Returns a deep copy of the \classpymatrix as a \pythcode{pandas.DataFrame}.}
\begin{lstlisting}[style=pythonstylefoot]
.to_pandas() -> pandas.core.frame.DataFrame
\end{lstlisting}
\generaldesc{This method returns a new, independent
	\pythcode{pandas.DataFrame} containing the same data as the
	original \pythcode{gretl.Matrix}. Changes made to the
	returned DataFrame do not affect the original matrix.}
